\section{A first nonprojective tensor-power module}
\label{sec:projectivity}

Retain the projective generator $P$ and the exact faithful evaluation
$\omega=\operatorname{Hom}_{\mathcal C}(P,-)$ from
Lemma~\ref{lem:finite-conventions}. We now select the rank of the elementary
abelian group for the higher Frobenius construction. The essential point is that exact
evaluation by a projective generator detects internal norm maps, even
though that evaluation does not preserve tensor products.

\subsection{Norms and independent group actions}

For an object $X$ with an action of a finite group $H$, denote its
invariants and coinvariants by $X^H$ and $X_H$.  These are, respectively,
the simultaneous kernel and cokernel of the maps $h-1$.  The endomorphism
$\sum_{h\in H}h$ factors through a morphism
\[
 \nu_H(X):X_H\longrightarrow X^H,
\]
called the norm map.  Exact $k$-linear functors preserve invariants,
coinvariants, and this map.

\begin{lemma}\label{lem:norm-projectivity}
Let $H$ be a finite $p$-group and $M$ a finite-dimensional $kH$-module.
Then $M$ is projective if and only if $\nu_H(M)$ is an isomorphism.
For an arbitrary finite group $G$, the norm map of every projective
$kG$-module is an isomorphism.
\end{lemma}

\begin{proof}
For the regular module $kG$, both coinvariants and invariants are
one-dimensional, generated by $[1]$ and $\sum_{g\in G}g$, respectively;
the norm takes the first generator to the second.  The assertion therefore
holds for free modules and their direct summands.

For the converse, $kH$ is a local self-injective algebra.  Its unique
simple module is the trivial module, and its socle is the line generated
by $N_H=\sum_{h\in H}h$.  The augmentation ideal is nilpotent by induction
on $|H|$: choose a central element $z$ of order $p$, use $(z-1)^p=0$,
and pass to $k[H/\langle z\rangle]$.  This proves locality.  The
nondegenerate associative pairing that takes $(a,b)$ to the coefficient
of $1$ in $ab$ identifies $kH$ with its vector-space dual as a module,
proving self-injectivity.  The socle is the space of left $H$-invariants
in the regular module, hence the stated line.
Every nonzero left ideal meets the socle.  Thus, if $N_Hm\ne0$, the map
$kH\to M$, $a\mapsto am$, is injective and splits by self-injectivity.
Removing free summands repeatedly gives $M\simeq(kH)^a\oplus M'$ with
$N_HM'=0$.  If $M'\ne0$, then $M'_H\ne0$ by Nakayama's Lemma, so its
zero norm map cannot be an isomorphism.  An isomorphic norm therefore
forces $M'=0$.
\end{proof}

\begin{lemma}\label{lem:internal-norm-products}
Let finite $p$-groups $H_1,\dots,H_s$ act on objects
$X_1,\dots,X_s$ of $\mathcal C$, respectively.  If each
$\omega(X_i)$ is projective over $kH_i$, then
$\omega(X_1\otimes\cdots\otimes X_s)$ is projective over
$k[H_1\times\cdots\times H_s]$, for the independent actions on the
tensor factors.  Its restriction to every subgroup of that product is
projective as well.
\end{lemma}

\begin{proof}
By Lemma~\ref{lem:norm-projectivity}, each evaluated norm
$\omega(\nu_{H_i}(X_i))$ is an isomorphism.  Exact faithfulness of
$\omega$ implies that each internal norm is an isomorphism.
Exactness of tensor products gives the natural identifications
\[
 (X\otimes Y)_{H\times K}\simeq X_H\otimes Y_K,
 \qquad
 (X\otimes Y)^{H\times K}\simeq X^H\otimes Y^K
\]
when $H$ acts on $X$ and $K$ acts on $Y$.  One obtains these identities
by taking kernels or cokernels for one group at a time.  Under them,
$\nu_{H\times K}(X\otimes Y)=\nu_H(X)\otimes\nu_K(Y)$.
Iteration proves that the product-group norm is an isomorphism.
Applying $\omega$ and Lemma~\ref{lem:norm-projectivity} proves the first
claim.  Restriction preserves projectivity because a finite group algebra
is free over the group algebra of any subgroup.
\end{proof}

\subsection{Existence and minimality}

Let $E_a=(\mathbb Z/p)^a$ act on $p^a$ positions by its regular action,
and hence act on $P^{\otimes p^a}$ through the symmetry of $\mathcal C$.
Define
\[
 M(a)=\omega(P^{\otimes p^a})\qquad(a\ge1).
\]

\begin{proposition}\label{prop:first-nonprojective}
There is a least integer $r\ge1$ for which $M(r)$ is not projective over
$kE_r$.  For this $r$, the restriction of $M(r)$ to every proper subgroup
of $E_r$ is projective.
\end{proposition}

\begin{proof}
Suppose first that every $M(a)$ is projective.  Fix $d\ge1$ and an
elementary abelian $p$-subgroup $H\subseteq\mathfrak S_d$.
On each $H$-orbit, the action is regular for a quotient $H/K_i$, because
$H$ is abelian.  Choose identifications $H/K_i\simeq E_{a_i}$, allowing
$a_i=0$ for singleton orbits.  Regrouping tensor factors identifies
$P^{\otimes d}$ with the tensor product of the corresponding orbit
powers.  Each evaluated orbit power is projective for its quotient
group, by the assumption when $a_i>0$, and trivially when $a_i=0$.
Lemma~\ref{lem:internal-norm-products} applies to the independent
product of the quotient groups.  The homomorphism
$H\to\prod_iH/K_i$ is injective, since the permutation action of $H$
is faithful.  Restriction along it shows that
$\omega(P^{\otimes d})$ is projective over $kH$.
Chouinard's projectivity theorem now implies that
$\omega(P^{\otimes d})$ is projective over $k\mathfrak S_d$
\cite{Chouinard1976}; see also
\cite[p.~20]{Benson2015}.

The tensor power of $P\twoheadrightarrow\mathbf1$ is an
$\mathfrak S_d$-equivariant epimorphism to the unit with its trivial
action.  Exact evaluation, followed by coinvariants, therefore gives an
epimorphism
\[
 \omega(P^{\otimes d})_{\mathfrak S_d}
 \twoheadrightarrow\omega(\mathbf1)\ne0.
\]
Lemma~\ref{lem:norm-projectivity} shows that the total symmetrizer
$b_d=\sum_{\sigma\in\mathfrak S_d}\sigma$ acts nontrivially on
$\omega(P^{\otimes d})$, and hence on $P^{\otimes d}$.

On the other hand, the action maps
\[
 k\mathfrak S_d\longrightarrow
 \operatorname{End}_{\mathcal C}(P^{\otimes d})
\]
cannot all be injective.  To see this without any dimension theory,
write $\ell(X)$ for the length of $X$ and put
$b=\max_i\ell(S_i\otimes P)$.  Exactness gives
$\ell(P^{\otimes d})\le\ell(P)b^{d-1}$.
Since the simples have endomorphism field $k$, induction on two
composition series gives
$\dim_k\operatorname{Hom}(X,Y)\le\ell(X)\ell(Y)$.
Thus $\dim_k\operatorname{End}(P^{\otimes d})$ grows at most
exponentially, whereas $\dim_k k\mathfrak S_d=d!$.

For clarity, the kernels form a tensor ideal in the permutation category
$\mathsf{Sym}$: its objects are the nonnegative integers, its
endomorphism algebra at $d$ is $k\mathfrak S_d$, other morphism spaces
are zero, and tensor product is addition with the usual block inclusions
of symmetric groups.  The symmetry of $\mathcal C$ gives a monoidal
functor $\mathsf{Sym}\to\mathcal C$, $d\mapsto P^{\otimes d}$, and
the kernels form its nonzero kernel ideal.  By
\cite[Lemma~4.10]{CoulembierEtingofOstrik2023}, every nonzero tensor
ideal in $\mathsf{Sym}$ contains some total symmetrizer $b_d$ in
positive characteristic. Its proof uses Kleshchev's faithfulness result
in the appendix to that paper \cite[Appendix~A, Proposition~A.1]{CoulembierEtingofOstrik2023}.
This contradicts the preceding nonvanishing.
The least $r$ therefore exists.

Finally, let $H<E_r$ be proper, of rank $a<r$.  Its action on the
regular $E_r$-set is a disjoint union of regular $H$-orbits.  Each orbit
power has projective evaluation by minimality when $a>0$, and trivially
when $a=0$.  Apply Lemma~\ref{lem:internal-norm-products} to the
independent orbit groups, then restrict to their diagonal copy of $H$.
This proves the last assertion, including $r=1$.
\end{proof}

The proposition supplies a nonprojective $E_r$-module whose restriction
to every proper subgroup is projective.  The next step uses its support
to choose the embedding $E_r\hookrightarrow\mathbb G_a$.
